\begin{scanpage}{0010}{vi}
This name is applied to many works on the diffraction of short waves in which the mathematical difficulties are bypassed by means of physical considerations.

It is clear that the physical theory of diffraction is a step forward in comparison with physical optics, which in general neglects the additional (edge) currents. The results obtained in this book show that with a given wavelength the physical theory of diffraction gives a better precision than physical optics, and with a given precision the physical theory of diffraction allows one to advance into the longer wave region and, in particular, to obtain a number of results which are of interest for radar where the ratios of the dimensions of the bodies to the wavelength do not reach such large values as in optics.

In addition, the physical theory of diffraction encompasses a number of interesting phenomena which are entirely foreign to physical optics. Thus, in a number of cases the additional currents give, not a small correction to the radiation field, but the main contribution to this field (see especially Chapters IV and V). If a plane wave is diffracted by a thin straight wire (a passive vibrator), then the additional current falls off very slowly as one goes further from the end of the wire. Therefore, the solution is obtained by summing the entire array of diffraction waves (secondary, tertiary, etc.) which successively arise as a consequence of the reflection of the currents from the ends of the wires. It has a resonance character. Thus, the problem of the scattering of the plane wave by a finite length wire which is a diffraction problem of a slightly unusual type is solved in Chapter VII. The resulting solution is applicable under the condition that the diameter of the wire is small in comparison with the wavelength and length of the wire, and the ratio of the length of the wire to the wavelength is arbitrary.

The final equations which are derived in this book and are used for calculations are not asymptotic in the strict sense of the word.
\end{scanpage}
